\subsection{状态契约、参数标定与验证边界}
\subsubsection{任务层转换契约与可达状态}
表中把路线图转换成可核对的前置条件、执行行为和后继状态。状态可达性不仅由视觉决定，也由运行模式决定：no-grab绕过抓取播放，no-release绕过终点放置，交互暂停增加等待恢复，无头失败直接结束。分析完整任务时必须固定这四类开关，否则“结束”对应不同物理结果。
\begin{longtable}{p{25mm}p{38mm}p{42mm}p{27mm}}\caption{任务主状态的转换契约。}\\\toprule 当前状态 & 事件/守卫 & 执行动作 & 后继状态\\\midrule\endfirsthead\toprule 当前状态 & 事件/守卫 & 执行动作 & 后继状态\\\midrule\endhead
等待 & S或无头延时结束 & 放行go & 可选抓取\\
抓取 & 动作播放返回成功 & 从尾帧读取六个锁臂值 & 走向A\\
走向A & A距离到阈值/近距推断 & 记录corner，调用left & 找B\\
找B & B满足新帧多数可见 & corner返回True & 走向B\\
走向B & B距离到阈值/近距推断 & 记录corner，调用right & 找C\\
找C & C满足新帧多数可见 & corner返回True & 走向C\\
走向C & C距离到阈值/近距推断 & 记录stop，按开关放置 & 任务返回\\
任意执行阶段 & 停止事件/异常 & 停止后续帧，记录原因 & 退出/清理\\\bottomrule\end{longtable}
抓取、保持和放置应视为独立物理谓词，但当前入口只对动作序列作控制，不独立验证所有物理谓词。可以定义规划任务成功为P抓取且P路线且P放置，线程正常返回仅是软件条件，不能与该合取等价。正式验收应使用连续视频或新增传感器给这些谓词赋值，附上对应日志时间。
\source{\code{marker_route.py}：Route.\_run、main及ROUTE。}

\subsubsection{单段决策的优先级及边界输入}
分支顺序决定相同观测怎样被解释。无相机新帧先等待；随后判断近距丢码；有效观测先判到达，再判方位；只有未到达且方位在容差内才前进。若距离已经小于阈值但方位很偏，程序仍优先认定到达，这与“必须先严格居中再到达”的控制律不同。报告不改源码，仅把该优先级作为需要现场核对的边界。
\begin{table}[H]\centering\small\begin{tabularx}{\linewidth}{p{28mm}YY}\toprule 输入情形 & 当前响应 & 验证重点\\\midrule
不足三个新帧 & 相机等待 & 断流不持续走/转\\
三帧仅一帧有目标 & 目标不可见分支 & 盲行上限及先前是否见过\\
距离达到，方位偏大 & 到达优先 & 路口空间与码布置是否合理\\
距离未到，方位超门限 & 单次纠偏 & 防来回摆与最大纠偏次数\\
距离未到，方位容差内 & 单步直行 & 负载下实际偏航和距离推进\\
近距最后观测后丢码 & lost\_close到达推断 & 遮挡造成误触发的风险\\\bottomrule\end{tabularx}\caption{单段决策的边界案例，不是新增控制器。}\end{table}
对每个边界案例应使用固定输入或真实日志回放验证动作选择，算法逻辑验证和物理试验分别记录。临界值附近的跳变并不必然是bug，因为离散动作本身采用阈值切换；重要的是单步幅度和迟滞是否与实际机体相容，以及守卫是否防止状态无限抖动。
\source{\code{marker_route.py}：obs判断、arrived\_blind与need分支。}

\subsubsection{相机资源重连、USB故障与视觉时效}
\texttt{robust\_\allowbreak camera}使用USB身份查找摄像头，避免设备号变化后固定video0失效，并处理占用相机的网页进程与断线重连。共享摄像头被两个程序分别打开时，能打开节点不代表可以持续取得新图；相机恢复后还要确认帧更新时间与标记检测线程确实推进。报告应区分设备存在、取流成功、新帧连续和识别有效四个层级。

USB断开可能与供电、线缆或大负载动作有关。当前材料只支持当次恢复和局部取流结果，不能凭重新插接一次就认定完整搬运期间不再掉线。建议把相机错误时间与动作类型、供电、电量、帧序号和等待时长对齐；若故障集中发生在转身，排查电源和线缆拉扯比单纯增大camera-wait更直接。
故障排查应从最早出现异常的层级向下追踪：有图像不等于有新帧，有新帧不等于目标有效，有效目标也不等于测距正确。按层保存证据可以避免将相机取流故障误判为路线决策错误。

\begin{table}[H]\centering\small\begin{tabularx}{\linewidth}{p{28mm}Y}\toprule 证据 & 可以支持的结论\\\midrule
设备节点可见 & 内核当前枚举到相机，不能证明能取图\\
连续新帧与序号 & 当前取流持续，不保证标记一定有效\\
标记多数可见 & 当前look观测成立，不保证后续转角可靠\\
完整负载录像+日志 & 可评价动作期间断流、等待与恢复效果\\\bottomrule\end{tabularx}\caption{视觉链故障证据的层级。}\end{table}

当取流已恢复但目标仍无效时，应继续检查标记尺寸、遮挡和检测输出；只有前置观测满足要求，路线状态转换才有可靠依据。按表中层级归档能够保留这一判断顺序。

\source{\code{robust_camera.py}；\code{Eyes.loop}；\code{Eyes.look}。}

\subsubsection{动作时间预算与感知控制带宽}
离散双足动作的控制频率由动作长度和观测等待共同决定。直行一次除了动作文件时长，还有\texttt{SETTLE\_\allowbreak S}=0.25 s、\texttt{WALK\_\allowbreak PAUSE\_\allowbreak S}=0.2 s和look取得三个新帧的时间；转角还包含前后各0.3 s。机械运动与感知并不是连续毫秒级闭环，因此用相机帧率当机器人纠偏带宽会严重高估响应。
\begin{equation}
T_{\mathrm{walkcycle}}\ge T_{\mathrm{gait}}+0.25+0.20+T_{\mathrm{look}},\qquad f_{\mathrm{control}}\le T_{\mathrm{walkcycle}}^{-1}.
\end{equation}
首帧桥接可能增加时长，断流和暂停会进一步延长。测量全程时间时应分解有效动作、正常等待、保护等待和人工恢复；较长总时间不必然表明视觉运算慢。优化方向先判断瓶颈：若多数时间用于大角度纠偏，调直抱物步态比提高图像线程帧率更有效；若look经常超时，则需要检查相机与计算资源。
\source{\code{marker_route.py}：WALK\_PAUSE\_S、CORNER\_PAUSE\_S、SETTLE\_S与act。}

\subsubsection{标定计划：从印刷尺寸到负载步态}
标定应分开相机几何、标记布局和身体运动。相机几何检查内参、畸变、实际分辨率以及14 cm黑边尺寸；布局记录A/B/C安装位置、高度、朝向与观察距离；运动测量抓取尾帧、负载质量、单步长度、左右偏航和转角分布。混合修改所有参数会使效果无法归因，也难以复现。
\begin{longtable}{p{28mm}p{45mm}p{59mm}}\caption{人形标定与数据归档建议。}\\\toprule 标定对象 & 记录量 & 验证方法\\\midrule\endfirsthead\toprule 标定对象 & 记录量 & 验证方法\\\midrule\endhead
标记打印 & 黑色方形实际边长 & 尺量多边，与\texttt{SIZE\_\allowbreak CM}一致\\
相机 & 内参、畸变、分辨率 & 多姿态图，独立距离检查\\
码布局 & 墙面坐标、高度与朝向 & 测绘图与相机视点照片\\
单步直行 & 负载、距离、偏航 & 多次独立起步，报告分布\\
单次转向 & 左/右角度与侧移 & 满载独立重复，外部量角\\
抓放衔接 & 首尾姿态、掉落 & 连续视频和动作文件哈希\\
相机恢复 & 断帧时长、恢复状态 & 人工断流与低速动作试验\\\bottomrule\end{longtable}
文档中的3.4 cm步长或10至15度偏航属于已有描述，不能视为新实验复测结果。报告引用这类数值要保留来源和条件，正式统计须使用同一设备、同一负载和同一动作版本，避免把不同阶段参数拼成一组“最终实测”。
\source{\code{maze_route/README.md}；\code{marker_route.py}运行参数。}

\subsubsection{课程要求与当前人形闭环的对应}
当前代码解决了基于固定标记的路线推进、离散纠偏与有界转角搜索；课程则同时要求抓取成功、全程抱持、不触障碍、左右转弯和放置。程序已有的状态守卫与这些要求并非一一充分对应。理解这个差异有助于设计实验，而不是削弱实现价值。
\begin{table}[H]\centering\small\begin{tabularx}{\linewidth}{p{25mm}YY}\toprule 课程要求 & 当前实现支撑 & 仍需的证据/闭环\\\midrule
抱起目标物 & \texttt{grab\_\allowbreak final\_\allowbreak v2}播放与锁臂 & 独立抓取确认或视频\\
全程保持抱物 & 命令手臂固定 & 实际滑落检测及掉落点恢复\\
左右各一次转向 & A-left/B-right及下一码获取 & 实测路线、角度、无干预\\
不触碰围挡 & 码距离及次数/推进保护 & 机体包络、测距或实测接触\\
终点放下 & C到达后put\_down & 放置成功和物体位置确认\\\bottomrule\end{tabularx}\caption{实现机制与任务验收条件的对应。}\end{table}
下一阶段应把真实失败归因到感知、动作或状态接口，再决定是否需要新增传感器。不能用合成路线成功覆盖物理持物与避碰的缺项，也不把软件存在的功能认定为未参与成员的实际工作。组内协作见第\ref{sec:division}节，技术结果关联具体文件和运行证据。
\source{课程任务手册；\code{maze_route/marker_route.py}与已有运行说明。}
